\documentclass[12pt,a4paper]{article}
\usepackage[margin=25mm, headheight=15pt, headsep=10pt]{geometry}
\usepackage{fontspec}
\usepackage[table]{xcolor} % Note: table option is required for rowcolors
\usepackage{titlesec}
\usepackage{tocloft}
\usepackage{fancyhdr}
\usepackage{booktabs}
\usepackage{tcolorbox}
\tcbuselibrary{skins, breakable}
\usepackage[colorlinks=true, linkcolor=emerald, urlcolor=emerald, bookmarks=true]{hyperref}

\definecolor{emerald}{RGB}{0,102,51}
\definecolor{darkred}{RGB}{139,0,0}
\definecolor{lightgray}{RGB}{245,245,245}

\usepackage[nil]{babel}
\babelprovide[import, main]{bengali}
\babelprovide[import]{arabic}
\babelfont{rm}[Renderer=HarfBuzz,Scale=1.0]{Noto Serif Bengali}
\babelfont{sf}[Renderer=HarfBuzz]{Noto Sans Bengali}
\babelfont{tt}[Renderer=HarfBuzz]{Noto Sans Bengali}
\babelfont[arabic]{rm}[Renderer=HarfBuzz,Scale=1.1]{Noto Naskh Arabic}

% Section styling
\titleformat{\section}{\Large\bfseries\color{emerald}}{\thesection.}{0.5em}{}
\titleformat{\subsection}{\large\bfseries\color{emerald!85!black}}{\thesubsection.}{0.5em}{}
\titleformat{\subsubsection}{\normalsize\bfseries\color{emerald!70!black}}{\thesubsubsection.}{0.5em}{}
\titlespacing*{\section}{0pt}{18pt}{8pt}

% Fancy Header/Footer setup
\pagestyle{fancy}
\fancyhf{} % clear all header and footer fields
\fancyhead[L]{\color{emerald}\small\textbf{\nouppercase{\leftmark}}}
\fancyfoot[L]{\scriptsize\color{black!50}{\lat{AI}}-সংকলিত, আলেম-যাচাইকৃত নয়}
\fancyfoot[C]{\color{emerald!70!black}\thepage}
\renewcommand{\headrulewidth}{0.4pt}
\renewcommand{\headrule}{\hbox to\headwidth{\color{emerald}\leaders\hrule height \headrulewidth\hfill}}
\renewcommand{\footrulewidth}{0pt}

% Custom boxes
% 1. Ayat/Hadith Box
\newtcolorbox{ayatbox}[1][]{
    enhanced, breakable, 
    colback=emerald!5, 
    colframe=emerald, 
    boxrule=0.5pt, 
    arc=3pt, 
    left=10pt, right=10pt, top=10pt, bottom=10pt,
    #1
}

% 2. Quote/Translation Box
\newtcolorbox{quotebox}[1][]{
    enhanced, breakable,
    frame hidden,
    colback=lightgray,
    borderline west={3pt}{0pt}{emerald},
    left=15pt, right=10pt, top=10pt, bottom=10pt,
    sharp corners,
    #1
}

% 3. AI-generated notice box
\newtcolorbox{warnbox}[1][]{
    enhanced, breakable,
    colback=red!4,
    colframe=darkred,
    boxrule=0.8pt,
    arc=3pt,
    left=10pt, right=10pt, top=10pt, bottom=10pt,
    #1
}

% Commands
\newcommand{\arab}[1]{\foreignlanguage{arabic}{#1}}
\newcommand{\kib}[1]{\textcolor{emerald}{\textbf{#1}}}

% Latin (English) fallback: Noto Serif Bengali has NO Latin glyphs
\newfontfamily\latfont[Renderer=HarfBuzz]{Noto Serif}
\newcommand{\lat}[1]{{\latfont #1}}

\setlength{\parskip}{5pt}
\setlength{\parindent}{0pt}

% Fix for missing bullet in Noto Serif Bengali
\renewcommand{\labelitemi}{$\bullet$}



\begin{document}
\begin{center}{\Large\bfseries\color{emerald} আল্লাহর রহমত ও তাওবা — \lat{06}-বেশ্যা-ও-তৃষ্ণার্ত-কুকুর}\end{center}
\vspace{1em}
\section*{ঘটনা ৬ — বেশ্যা ও তৃষ্ণার্ত কুকুর: একটি ভালো কাজে গোটা জীবনের ক্ষমা}

\begin{warnbox}
 \textbf{গুরুত্বপূর্ণ নোটিশ (\lat{Important} \lat{Notice}):} এই কন্টেন্টটি \lat{AI} (কৃত্রিম বুদ্ধিমত্তা) দিয়ে প্রস্তুত ও সংকলিত। \lat{AI} কঠোর সূত্র-মান নীতি মেনে শুধু নির্ভরযোগ্য উৎস থেকে তথ্য সংগ্রহ করেছে — কেবল সহীহ/হাসান হাদিস ও সহীহ সনদের আসার রাখা হয়েছে, দুর্বল সনদ বাদ দেওয়া হয়েছে। তবে এটি কোনো আলেম বা মুহাদ্দিস কর্তৃক যাচাইকৃত নয়।
\end{warnbox}

\begin{quote}
\textbf{সম্পর্কিত ফাইল:} হাদিস ফাইল — পশুর প্রতি দয়া; ব্যবহারিক গাইড — ভালো কাজ পাপ মুছে দেয়।
\end{quote}

\medskip\hrule\medskip

\subsection*{১. সূত্র কার্ড (\lat{Source} \lat{Card})}

\begin{table}[h]\centering\small
\begin{tabular}{ll}
বিষয় & বিবরণ \\
\hline
\textbf{ঘটনা} & এক পাপী নারী (বেশ্যা) একটি তৃষ্ণার্ত কুকুরকে পানি পান করাল — আল্লাহ তাকে ক্ষমা করলেন \\
\textbf{বর্ণনাকারী} & আবু হুরাইরা (রাঃ) \\
\textbf{গ্রন্থ ও নম্বর} & সহীহ বুখারি ৩৩২১; সহীহ মুসলিম ২২৪৫ \\
\textbf{অধ্যায়} & বুখারি: কিতাবু বাদয়িল খালক; মুসলিম: কিতাবুস সালাম — \lat{“}পশুকে পানি ও খাবার দেওয়ার ফজিলত\lat{”} \\
\textbf{সনদের মান} & \textbf{সহীহ} (বুখারি-মুসলিম) \\
\textbf{স্থান ও সময়} & বনি ইসরাঈলের যুগের ঘটনা — নবীজি (সা.) বর্ণনা করেছেন \\
\hline
\end{tabular}
\end{table}

\medskip\hrule\medskip

\subsection*{২. সংক্ষিপ্ত পটভূমি (\lat{Background})}

এটি ইসলামের সবচেয়ে বিখ্যাত ক্ষমার গল্প — পাপের আকার যত বড়ই হোক, রহমতের দরজা তার চেয়ে বড়। নারীটি পাপে নিমগ্ন ছিল — যিনা তার জীবিকা ও জীবনের অংশ ছিল; তবু একটি খাঁটি, নিঃস্বার্থ ভালো কাজ — এক প্রাণের তৃষ্ণা মেটানো — তার পুরো হিসাব বদলে দিল। অধ্যায়ের নামটি খেয়াল করুন — \textbf{\lat{“}পশুকে পানি দেওয়ার ফজিলত\lat{”}} — অর্থ: ভালো কাজের মূল্য তার আকারে নয়, তার ইখলাসে।

\medskip\hrule\medskip

\subsection*{৩. পূর্ণ ঘটনার বিশুদ্ধ বর্ণনা (\lat{The} \lat{Full} \lat{Narration})}

\begin{quote}
\textbf{\lat{“}এক বেশ্যা (পাপী নারী) একটি কুকুরকে কূপের পাশে দেখল — গরমের দিনে সে তৃষ্ণায় জিভ বের করে ছিল। সে তার জুতো খুলল, তা দিয়ে (দড়ি বেঁধে) পানি তুলল, কুকুরটিকে পান করাল। এ কারণে আল্লাহ তাকে ক্ষমা করে দিলেন।\lat{”}}
\end{quote}

\begin{quote}
মুসলিমের বর্ণনায়: \textbf{\lat{“}এক নারী — বেশ্যা — একটি কুকুরকে দেখল গরমের দিনে কূপের চারপাশে চক্কর দিচ্ছে, তৃষ্ণায় জিভ বের হয়েছে। সে তার জুতোর মাধ্যমে তার জন্য পানি তুলল — তাকে ক্ষমা করা হলো।\lat{”}}
\end{quote}

\begin{quote}
সহীহ বুখারির অপর বর্ণনায় (২৩৬৩) — \textbf{\lat{“}এক ব্যক্তি পথিমধ্যে এক কুকুরকে তৃষ্ণায় জিভ বের করতে দেখল, বলল: এই কুকুরটি তৃষ্ণায় এমন অবস্থায়, যেমন আমি ছিলাম। সে কূপে নামল, জুতো ভরে পানি তুলল, মুখ দিয়ে ধরে কুকুরটিকে পান করাল। আল্লাহ তার শোকর কবুল করলেন — তাকে ক্ষমা করলেন।\lat{”}}
\end{quote}

\medskip\hrule\medskip

\subsection*{৪. শব্দের অর্থ (\lat{Key} \lat{Words})}

\begin{table}[h]\centering\small
\begin{tabular}{ll}
শব্দ & অর্থ \\
\hline
\textbf{বাগিয়্যুন (\arab{بغيّ})} & বেশ্যা — যিনা-পেশায় লিপ্ত নারী \\
\textbf{ইয়ুতীফু বিইরিন (\arab{يطيف} \arab{ببئر})} & কূপের চারপাশে চক্কর দেওয়া \\
\textbf{আদলা'আ লিসানাহু} & জিভ বের করা — চরম তৃষ্ণা \\
\textbf{মুক্বুহা (\arab{موقها})} & তার জুতা \\
\textbf{ফাগুফিরা লাহা (\arab{فغفر} \arab{لها})} & তাই তাকে ক্ষমা করা হলো \\
\hline
\end{tabular}
\end{table}

\medskip\hrule\medskip

\subsection*{৫. সংশ্লিষ্ট দলিল ও রেফারেন্স}

\begin{quote}
\textbf{\lat{“}নিশ্চয়ই আল্লাহ সমস্ত গুনাহ ক্ষমা করে দেন।\lat{”}} — সূরা আয-যুমার ৩৯:৫৩
\end{quote}

\begin{quote}
\textbf{\lat{“}নিশ্চয়ই ভালো কাজগুলো মন্দ কাজগুলোকে মুছে দেয়।\lat{”}} — সূরা হুদ ১১:১১৪
\end{quote}

\begin{quote}
\textbf{\lat{“}মন্দ কাজের পর ভালো কাজ করো — তা তা মুছে দেবে।\lat{”}} — তিরমিযী ১৯৮৭ (হাসান-সহীহ)
\end{quote}

\subsubsection*{মূল শিক্ষা (দলিলভিত্তিক)}

\begin{enumerate}
\item \textbf{পাপের হিসাব রহমতের সামনে ছোট:} জীবনের যিনা — বড় পাপ — একটি ইখলাসের ভালো কাজে ক্ষমা হয়ে গেল; আল্লাহর রহমত পাপের আকার মাপে না, বান্দার ফেরার ইখলাস মাপে।
\item \textbf{ভালো কাজের মূল্য:} যেখানে দাঁড়িয়ে আছ — সেখানেই একটি ভালো কাজ রহমতের দরজা খুলে দিতে পারে (তিরমিযী ১৯৮৭) — পাপীকে \lat{“}ভালো হওয়ার আগে\lat{”} ভালো কাজ করতে বাধা নেই।
\item \textbf{সৃষ্টির প্রতি দয়া:} পশুর প্রতি দয়াও আল্লাহর কাছে মূল্যবান — রহমত পাওয়ার পথ সব দিকে খোলা।
\item \textbf{সতর্কতা:} হাদিস পাপকে ক্ষমা করার গল্প — পাপে থাকার অনুমতি নয়; এটি পাপের পর ফেরার পথ দেখায়, পাপ চালিয়ে যাওয়ার নয় (ব্যবহারিক গাইড ৩ নং দেখুন)।
\end{enumerate}
\medskip\hrule\medskip

\subsection*{৬. রেফারেন্স}

\begin{itemize}
\item সহীহ বুখারি ৩৩২১ (কিতাবু বাদয়িল খালক), ২৩৬৩ (কিতাবুল মুসাকাত)
\item সহীহ মুসলিম ২২৪৫ (কিতাবুস সালাম)
\item কুরআন: যুমার ৩৯:৫৩, হুদ ১১:১১৪
\item ব্যাখ্যা: ফাতহুল বারি, শারহু নববী
\end{itemize}

\end{document}
